\section*{Abstract}
		The T0 theory achieves complete parameter freedom: Only the geometric parameter $\xi = \frac{4}{3} \times 10^{-4}$ is fundamental. All physical constants are derived either from $\xi$ or represent unit definitions. For absolute SI values the derivation chain needs one mass or energy anchor; the speed of light is fixed by the definition of the metre and is not determined from $\xi$. This document provides the complete derivation chain including the gravitational constant $G$, the Planck length $l_P$, and the Boltzmann constant $k_B$. The SI reform 2019 unwittingly implemented the unique calibration consistent with this geometric foundation.

	
	
	
	
	\section{The Geometric Foundation}
	
	\subsection{Single Fundamental Parameter}
	
	\begin{equation}
		\boxed{\xi = \frac{4}{3} \times 10^{-4}}
	\end{equation}
	
	This geometric ratio encodes the fundamental structure of three-dimensional space. All physical quantities emerge as derivable consequences. (See \cite{T0_xi_origin_En} for the origin of $\xi$.)
	
	\subsection{Complete Derivation Framework}
	
	Detailed mathematical derivations are available at:
	
	\begin{center}
		
	\end{center}
	
	\section{Derivation of the Gravitational Constant from $\xi$}
	
	\subsection{The Fundamental T0 Gravitational Relationship}
	
	\begin{derivation}
		\textbf{Starting point of T0 gravity theory:}
		
		The T0 theory postulates a fundamental geometric relationship between the characteristic length parameter $\xi$ and the gravitational constant:
		
		\begin{equation}
			\xi = 2\sqrt{G \cdot m_{\text{char}}}
			\label{eq:t0_fundamental}
		\end{equation}
		
		where $m_{\text{char}}$ represents a characteristic mass of the theory. (Detailed derivation in \cite{T0_gravitational_constant_En}.)
		
		\textbf{Physical interpretation:}
		\begin{itemize}
			\item $\xi$ encodes the geometric structure of space
			\item $G$ describes the coupling between geometry and matter
			\item $m_{\text{char}}$ sets the characteristic mass scale
		\end{itemize}
	\end{derivation}
	
	\subsection{Solution for the Gravitational Constant}
	
	Solving equation \eqref{eq:t0_fundamental} for $G$:
	
	\begin{equation}
		\boxed{G = \frac{\xi^2}{4 m_{\text{char}}}}
		\label{eq:g_fundamental}
	\end{equation}
	
	This is the fundamental T0 relationship for the gravitational constant in natural units. (Further details in \cite{gravitational_constant_En}.)
	
	\subsection{Choice of Characteristic Mass}
	
	\begin{insight}
		\textbf{The electron mass follows from the $\xi$ mass ladder:}
		
		The T0 theory uses the electron mass as the characteristic scale:
		\begin{equation}
			m_{\text{char}} = m_e = 0{,}511 \text{ MeV}
			\label{eq:characteristic_mass}
		\end{equation}
		
		\textbf{Critical point:} The electron mass itself is not an independent parameter but follows from the T0 mass ladder (Doc.~006):
		\begin{equation}
			m_e = \frac{4}{3}\,\xi^{3/2}\,v
		\end{equation}
		
		The ladder needs $v$ as a scale; the chain $\xi \to m_e$ therefore does not work without a mass or energy anchor. (See \cite{T0_particle_masses_En} for mass derivations.)
		
		Therefore, the derivation chain $\xi \to m_e \to G \to l_P$ depends on $\xi$ as the single fundamental parameter and on one anchor for the scale.
	\end{insight}

	\subsection{Dimensional Analysis in Natural Units}
	
	\begin{derivation}
		\textbf{Dimensional check in natural units ($\hbar = c = 1$):}
		
		In natural units:
		\begin{align}
			[M] &= [E] \quad \text{(from } E = mc^2 \text{ with } c = 1\text{)} \\
			[L] &= [E^{-1}] \quad \text{(from } \lambda = \hbar/p \text{ with } \hbar = 1\text{)} \\
			[T] &= [E^{-1}] \quad \text{(from } \omega = E/\hbar \text{ with } \hbar = 1\text{)}
		\end{align}
		
		The gravitational constant has dimension:
		\begin{equation}
			[G] = [M^{-1}L^3T^{-2}] = [E^{-1}][E^{-3}][E^2] = [E^{-2}]
		\end{equation}
		
		Checking equation \eqref{eq:g_fundamental}:
		\begin{equation}
			[G] = \frac{[\xi^2]}{[m_e]} = \frac{[1]}{[E]} = [E^{-1}] \neq [E^{-2}]
		\end{equation}
		
		This shows that additional factors are required for dimensional correctness. (See \cite{natural_units_systematics_En} for unit systematics.)
	\end{derivation}
	
	\subsection{Complete Formula with Conversion Factors}
	
	\begin{keyresult}
		\textbf{Complete gravitational constant formula:}
		
		\begin{equation}
			\boxed{G_{\text{SI}} = \frac{\xi_0^2}{4 m_e} \times C_{\text{conv}} \times K_{\text{frak}}}
			\label{eq:G_complete}
		\end{equation}
		
		where:
		\begin{itemize}
			\item $\xi_0 = 1{,}333 \times 10^{-4}$ (geometric parameter)
			\item $m_e = 0{,}511$ MeV (electron mass, anchor of the chain)
			\item $C_{\text{conv}} = 7{,}783 \times 10^{-3}$ (SI conversion of the single-anchor chain, Doc.~383)
			\item $K_{\text{frak}} = 0{,}986$ (fractal quantum spacetime correction) (See \cite{fractal_correction_derivation_En}.)
		\end{itemize}
		
		\textbf{Result:}
		\begin{equation}
			G_{\text{SI}} = 6{,}674 \times 10^{-11} \text{ m}^3/(\text{kg}\cdot\text{s}^2)
		\end{equation}
		
		with $+0.003\%$ deviation from the CODATA-2018 value ($6.67430 \times 10^{-11}$) for $\xi = 4/30000$; with the rounded $\xi_0 = 1.333 \times 10^{-4}$ the deviation is $-0.05\%$.
		
		Here $C_{\text{conv}}$ is the SI conversion of the single-anchor chain (Doc.~383); its numerical value $7.783\times10^{-3}$ additionally contains the small residual of this chain (chain value $7.58\times10^{-3}$). The form of $G$ and its order of magnitude follow from $\xi$ and the one anchor (numerically checked in Doc.~383); with the anchor $m_e$, $G$ lies at $-2.6\,\%$, with $v$ at $-0.46\,\%$. The agreement to $0.003\,\%$ arises only with this residual and does not count as an independent precision result.
	\end{keyresult}
	
	\section{Derivation of the Planck Length from $G$ and $\xi$}
	
	\subsection{The Planck Length as Fundamental Reference}
	
	\begin{derivation}
		\textbf{Definition of the Planck length:}
		
		In standard physics, the Planck length is defined as:
		\begin{equation}
			l_P = \sqrt{\frac{\hbar G}{c^3}}
			\label{eq:planck_length_standard}
		\end{equation}
		
		In natural units ($\hbar = c = 1$) this simplifies to:
		\begin{equation}
			\boxed{l_P = \sqrt{G} = 1 \quad \text{(natural units)}}
			\label{eq:planck_natural}
		\end{equation}
		
		\textbf{Physical meaning:} The Planck length represents the characteristic scale of quantum gravitational effects and serves as the natural length unit in theories combining quantum mechanics and general relativity. (See \cite{T0_natural_SI_En} for natural and SI units.)
	\end{derivation}
	
	\subsection{T0 Derivation: Planck Length via $G_{\text{SI}}$}
	
	\begin{keyresult}
		\textbf{Complete derivation chain:}
		
		Since $G$ is derived from $\xi$ via equation \eqref{eq:g_fundamental}:
		\begin{equation}
			G = \frac{\xi^2}{4 m_e}
		\end{equation}
		
		the Planck length follows via the complete $G_{\text{SI}}$ from Eq.~\eqref{eq:G_complete}, which contains $C_{\text{conv}}$ and $K_{\text{frak}}$:
		\begin{equation}
			l_P = \sqrt{\frac{\hbar\,G_{\text{SI}}}{c^3}}, \qquad G_{\text{SI}} = \frac{\xi^2}{4 m_e}\,C_{\text{conv}}\,K_{\text{frak}}
		\end{equation}

		\textbf{Result in SI units:}
		\begin{equation}
			\boxed{l_P = 1{,}616 \times 10^{-35} \text{ m}}
		\end{equation}

	\end{keyresult}
	
	\subsection{The Characteristic T0 Length Scale}
	
	\begin{insight}
		\textbf{Connection between $r_0$ and the fundamental energy scale $E_0$:}
		
		The characteristic T0 length $r_0$ for an energy $E$ is defined as:
		\begin{equation}
			r_0(E) = 2GE
		\end{equation}
		
		For the fundamental energy scale $E_0 = \sqrt{m_e \cdot m_\mu}$:
		\begin{equation}
			r_0(E_0) = \frac{\hbar c}{E_0} \approx 2.7 \times 10^{-14} \text{ m}
		\end{equation}
		
		The minimal sub-Planck length scale is:
		\begin{equation}
			\boxed{L_0 = \xi \cdot l_P = \frac{4}{3} \times 10^{-4} \times 1{,}616 \times 10^{-35} \text{ m} = 2{,}155 \times 10^{-39} \text{ m}}
		\end{equation}
		
		\textbf{Fundamental relationship:} In natural units, for any energy $E$:
		\begin{equation}
			r_0(E) = \frac{1}{E} \quad \text{(in natural units with } c = \hbar = 1\text{)}
		\end{equation}
		
		where the time-energy duality $r_0(E) \leftrightarrow E$ defines the characteristic scale. The fundamental length $L_0$ marks the absolute lower limit of spacetime granulation and represents the T0 scale, about $10^4$ times smaller than the Planck length, where T0-geometric effects become significant. (See \cite{T0_energy_En} for energy scales.)
	\end{insight}
	
	\subsection{The Decisive Convergence: Why T0 and SI Agree}
	
	\begin{historical}
		\textbf{Two paths to the same Planck length:}
		
		There are two ways to determine the Planck length:
		
		\textbf{Path 1: SI-based (experimental):}
		\begin{equation}
			l_P^{\text{SI}} = \sqrt{\frac{\hbar G_{\text{measured}}}{c^3}} = 1{,}616 \times 10^{-35} \text{ m}
		\end{equation}
		
		This uses the experimentally measured gravitational constant $G_{\text{measured}} = 6{,}674 \times 10^{-11}$ m$^3$/(kg$\cdot$s$^2$) from CODATA.
		
		\textbf{Path 2: T0-based (geometry and one anchor):}
		\begin{align}
			m_e &= \frac{4}{3}\,\xi^{3/2}\,v \quad \text{(from } \xi \text{ and the anchor } v\text{)} \\
			G &= \frac{\xi^2}{4m_e} \times C_{\text{conv}} \times K_{\text{frak}} \quad \text{(from } \xi \text{ and } m_e\text{)} \\
			l_P^{\text{T0}} &= \sqrt{\frac{\hbar G}{c^3}} \quad \text{(with } G = G_{\text{SI}} \text{ from Eq.~\eqref{eq:G_complete})}
		\end{align}
		
		\textbf{Conversion to SI units:}
		\begin{equation}
			l_P^{\text{T0}} = \sqrt{\frac{\hbar}{c^3} \cdot \frac{\xi^2}{4m_e} \, C_{\text{conv}} \, K_{\text{frak}}} = \sqrt{\frac{\hbar \cdot 6.6745 \times 10^{-11}}{c^3}} \text{ m}
		\end{equation}
		
		\textbf{Result:} $l_P^{\text{T0}} = 1.616 \times 10^{-35}$ m
		
		\textbf{The convergence:}
		\begin{equation}
			\boxed{l_P^{\text{SI}} = l_P^{\text{T0}} \quad \text{with } +0.002\% \text{ deviation}}
		\end{equation}
		As for $G$, this precision arises only with the residual in $C_{\text{conv}}$ and does not count as an independent precision result; the deviation is half that of $G$.
		
	\end{historical}
	
	\begin{warning}
		\textbf{Why this agreement is not coincidental:}
		
		The agreement between the SI-derived and T0-derived Planck length in form and order of magnitude, which follow from $\xi$ and one anchor, reveals a profound truth:
		
		\begin{enumerate}
			\item The SI reform 2019 unwittingly calibrated itself to geometric reality
			
			\item Sommerfeld's calibration in 1916 to $\alpha \approx 1/137$ was not arbitrary -- it reflected the fundamental geometric value $\alpha = \xi \cdot E_0^2$ (See \cite{T0_fine_structure_En}.)
			
			\item The experimental measurement of $G$ does not determine an arbitrary constant -- it measures the geometric structure encoded in $\xi$
			
			\item \textbf{The conversion factor:} The factor $\frac{\hbar c}{1 \text{ MeV}} = 1{,}973 \times 10^{-13}$ m converts lengths from natural units ($r_0 = 1/E$) into metres; the agreement of the Planck length itself follows from the complete $G_{\text{SI}}$ (Eq.~\eqref{eq:G_complete}), not from $S_{T0}$.
			
			\item Both paths describe the same underlying geometric reality: \textbf{the universe is pure $\xi$-geometry}
		\end{enumerate}
		
		The SI constants ($c$, $\hbar$, $e$, $k_B$) define \emph{how we measure}, but the \emph{relationships between measurable quantities} are determined by $\xi$-geometry. Therefore, the SI reform 2019, by fixing these unit-defining constants, unwittingly implemented the unique calibration consistent with T0 theory.
	\end{warning}
	
	\section{The Role of the Conversion Factor}
	
	\subsection{The Scaling Factor $S_{T0} = 1$ MeV/$c^2$}
	
	\begin{keyresult}
		\textbf{$S_{T0} = 1$ MeV/$c^2$ holds by construction:}
		
		The mass scaling factor is
		\begin{equation}
			\boxed{S_{T0} = \frac{m_e^{\text{SI}}}{m_e^{\text{T0}}} = 1 \text{ MeV}/c^2}
		\end{equation}
		
		with $m_e^{\text{T0}} = 0.511$. This value arises by construction, because $0.511$ is the electron mass in MeV (cf.\ Doc.~014); it is not a prediction of the MeV scale. The Planck length agrees through $G_{\text{SI}}$ and $C_{\text{conv}}$, not through $S_{T0}$. Where $\alpha$ hides in the conversion and what role the reference energy $1$~MeV plays is shown in Doc.~386: the agreement of $\alpha$ requires a reference energy of $0.99992$~MeV, whose geometric meaning is open. (See \cite{parameter_derivation_En} for parameter derivation.)
	\end{keyresult}
	
	\subsection{The Conversion Chain}
	
	\begin{derivation}
		\textbf{From natural units to SI units:}
		
		The conversion factor between natural T0 units and SI units is:
		\begin{equation}
			\text{Conversion factor} = \frac{\hbar c}{S_{T0}} = \frac{\hbar c}{1 \text{ MeV}} = 1{,}973 \times 10^{-13} \text{ m}
		\end{equation}
		
		For the Planck length (with $G_{\text{SI}}$ from Eq.~\eqref{eq:G_complete}):
		\begin{align}
			l_P^{\text{SI}} &= \sqrt{\frac{\hbar G_{\text{SI}}}{c^3}} \\
			&= \sqrt{\frac{1.0546 \times 10^{-34} \times 6.6745 \times 10^{-11}}{(2.9979 \times 10^{8})^3}} \text{ m} \\
			&= 1.616 \times 10^{-35} \text{ m} \quad \checkmark
		\end{align}
		
		\textbf{The geometric lock:} The agreement of the T0-derived Planck length with the SI-measured value is established through $G_{\text{SI}}$, which contains $\xi$, $m_e$ (in MeV, i.e.\ with $S_{T0} = 1$ MeV/$c^2$) and $C_{\text{conv}}$.
	\end{derivation}
	
	\subsection{The Triple Consistency}
	
	\begin{insight}
		\textbf{Three independent measurements lock together:}
		
		The system is overdetermined by three independent experimental values:
		\begin{enumerate}
			\item Fine structure constant: $\alpha = 1/137{,}035999084$ (determined from $a_e$ and atom-recoil measurements $h/m$ on Cs and Rb) (See \cite{fine_structure_constant_En}.)
			\item Gravitational constant: $G = 6{,}674 \times 10^{-11}$ m$^3$/(kg$\cdot$s$^2$) (Cavendish-type experiments)
			\item Planck length: $l_P = 1{,}616 \times 10^{-35}$ m (derived from $G$, $\hbar$, $c$)
		\end{enumerate}
		
		The T0 theory connects all three through $\xi$ and one mass or energy anchor: $\alpha$ via $\alpha = \xi \cdot E_0^2$, $G$ and $l_P$ via $G_{\text{SI}}$ with $C_{\text{conv}}$ and $K_{\text{frak}}$. The scaling factor $S_{T0} = 1$ MeV/$c^2$ is not an additional boundary condition here but holds by construction. The form and order of magnitude of $G$ and $l_P$ follow from $\xi$ and the anchor; the agreement to $0.003\,\%$ and $0.002\,\%$ respectively arises only with the residual in $C_{\text{conv}}$ and does not count as an independent precision result.
	\end{insight}
	
	\section{The Speed of Light: Geometric or Conventional?}
	
	\subsection{The Dual Nature of $c$}
	
	\begin{derivation}
		\textbf{Understanding the role of the speed of light:}
		
		The speed of light has a subtle dual character requiring careful analysis:
		
		\textbf{Perspective 1: As a dimensional convention}
		
		In natural units, setting $c = 1$ is purely conventional:
		\begin{equation}
			[L] = [T] \quad \text{(space and time have the same dimension)}
		\end{equation}
		
		This is analogous to saying 1 hour equals 60 minutes -- it's a choice of measurement units, not physics. (See \cite{unit_conventions_c_speed_En}.)
		
		\textbf{Perspective 2: As a ratio of the Planck scales}
		
		Via the Planck scales, $c$ can be written as a ratio:
		\begin{align}
			l_P &= \sqrt{\hbar G/c^3}, \quad G = \frac{\xi^2}{4m_e}\,C_{\text{conv}}\,K_{\text{frak}} \\
			t_P &= \frac{l_P}{c} = \frac{l_P}{1} \quad \text{(in natural units)}
		\end{align}
		
		Since $t_P$ is defined as $l_P/c$, $c = l_P/t_P$ is an identity and not a determination of $c$ by $\xi$.
	\end{derivation}
	
	\subsection{The SI Value is Fixed by Definition}
	
	\begin{keyresult}
		\textbf{Why $c = 299\,792\,458$ m/s exactly:}
		
		Since 1983, $c = 299\,792\,458$ m/s has been fixed by the definition of the metre; the value was chosen to match centuries of measurements. The ratio of the Planck scales
		\begin{equation}
			c^{\text{SI}} = \frac{l_P^{\text{SI}}}{t_P^{\text{SI}}} = \frac{1{,}616 \times 10^{-35} \text{ m}}{5{,}391 \times 10^{-44} \text{ s}}
		\end{equation}
		reproduces this value because $t_P = l_P/c$ by definition; it is an identity and not a determination of $c$ by $\xi$.
	\end{keyresult}
	
	\subsection{The Meter is Defined by $c$}
	
	\begin{insight}
		\textbf{The circular calibration loop:}
		
		There is a beautiful circularity in the SI-2019 system:
		
		\begin{enumerate}
			\item The meter is \emph{defined} as the distance light travels in $1/299\,792\,458$ seconds
			\item But the number $299\,792\,458$ was chosen to match experimental measurements
			\item $c = l_P/t_P$ is an identity because $t_P = l_P/c$; via the Planck scales, $c$ is not determined by $\xi$
		\end{enumerate}
		
		\textbf{Conclusion:} We use $c$ to \emph{define} the meter; the numerical value of $c$ is thus a fixing of the unit of length and does not follow from $\xi$. (See \cite{circularity_constants_En} for circularity of constants.)
	\end{insight}
	
	\section{Derivation of the Boltzmann Constant}
	
	\subsection{The Temperature Problem in Natural Units}
	
	\begin{warning}
		\textbf{The Boltzmann constant is NOT fundamental:}
		
		In natural units, where energy is the fundamental dimension, temperature is just another energy scale. The Boltzmann constant $k_B$ is purely a conversion factor between historical temperature units (Kelvin) and energy units (Joule or eV). (See \cite{temperature_units_CMB_En} for temperature units.)
	\end{warning}
	
	\subsection{Definition in the SI System}
	
	\begin{derivation}
		\textbf{The SI reform 2019 definition:}
		
		Since May 20, 2019, the Boltzmann constant has been fixed by definition:
		\begin{equation}
			\boxed{k_B = 1{,}380649 \times 10^{-23} \text{ J/K}}
			\label{eq:kb_si}
		\end{equation}
		
		This defines the Kelvin scale in terms of energy:
		\begin{equation}
			k_B \cdot 1 \text{ K} = 1.380649 \times 10^{-23} \text{ J}
		\end{equation}
	\end{derivation}
	
	\subsection{Relationship to Fundamental Constants}
	
	\begin{keyresult}
		\textbf{Boltzmann constant from gas constant:}
		
		The Boltzmann constant is defined by the Avogadro number:
		\begin{equation}
			k_B = \frac{R}{N_A}
		\end{equation}
		
		where:
		\begin{itemize}
			\item $R = 8{,}314462618$ J/(mol$\cdot$K) (ideal gas constant)
			\item $N_A = 6{,}02214076 \times 10^{23}$ mol$^{-1}$ (Avogadro constant, fixed since 2019)
		\end{itemize}
		
		\textbf{Result:}
		\begin{equation}
			k_B = \frac{8{,}314462618}{6{,}02214076 \times 10^{23}} = 1{,}380649 \times 10^{-23} \text{ J/K}
		\end{equation}
	\end{keyresult}
	
	\subsection{T0 Perspective on Temperature}
	
	\begin{insight}
		\textbf{Temperature as an energy scale in T0 theory:}
		
		In T0 theory, temperature is naturally expressed as energy:
		\begin{equation}
			T_{\text{natural}} = k_B T_{\text{Kelvin}}
		\end{equation}
		
		For example, the CMB temperature:
		\begin{align}
			T_{\text{CMB}} &= 2{,}725 \text{ K} \\
			T_{\text{CMB}}^{\text{natural}} &= k_B \times 2{,}725 \text{ K} = 2{,}35 \times 10^{-4} \text{ eV}
		\end{align}
		
		\textbf{Core message:} $k_B$ is not derived from $\xi$ because it represents a historical convention for temperature measurement, not a physical property of spacetime geometry.

		\textbf{Distinction $\xi_{\text{FFGFT}}$ vs. $\xi_{\text{UIFT}}$:}
		In the IPI discussion, $\xi_{\text{UIFT}} = k_B T_{\text{CMB}} \ln 2 / (m_e c^2) \approx 3.19 \times 10^{-10}$ is treated as an independent parameter.
		$\xi_{\text{FFGFT}} = 4/30000 \approx 1.333 \times 10^{-4}$ and $\xi_{\text{UIFT}}$ are \textbf{not the same quantity} — their ratio is:
		\begin{equation}
			\frac{\xi_{\text{FFGFT}}}{\xi_{\text{UIFT}}} = \frac{\xi\, m_e c^2}{k_B T_{\text{CMB}} \ln 2} \approx 418\,521
		\end{equation}
		Replacing $k_B T_{\text{CMB}}$ by the FFGFT value $(16/9)\,\xi$~eV gives $\frac{m_e c^2/\text{eV}}{(16/9)\ln 2} \approx 414\,684$.
		The factor follows from the SI conversion: $T_{\text{CMB}}[\text{eV}] = k_B \cdot 2.72548\,\text{K} = 2.3486 \times 10^{-4}\,\text{eV}$, while FFGFT to first order gives $(16/9)\,\xi = 2.370 \times 10^{-4}\,\text{eV}$ ($\Delta = 0.93\,\%$).
	\end{insight}

	\subsection{Static and Dynamic Unit of Information}

	The factor $\xi_{\text{FFGFT}}/\xi_{\text{UIFT}} \approx 418\,521$ (with the FFGFT value for $k_B T_{\text{CMB}}$: $414\,684$) is more than a
	conversion: it separates two fundamentally different notions of ``information''.
	$\xi_{\text{UIFT}}$ is defined through $k_B T_{\text{CMB}} \ln 2$ -- i.e.\ through
	an \emph{energy per bit}, hence through a temperature. $\xi_{\text{FFGFT}}$ is
	dimensionless and reference-free. This is exactly where the question arises of how
	much energy corresponds to the smallest unit of information.

	\begin{insight}
		\textbf{There is no universal energy per bit.}
		The energy of a unit of information is not a fixed quantity but a projection
		that requires a scale:
		\begin{align}
			\text{thermal (Landauer):}\quad & E = k_B T \ln 2 \notag \\
			& \text{(temperature-dependent)} \\[4pt]
			\text{spatial (energy hierarchy):}\quad & E_N = \frac{\hbar c}{L_N} \notag \\
			& \text{(length-dependent)}
		\end{align}
		Both expressions carry the embedding constants ($k_B$, $\hbar$, $c$) and an
		external scale ($T$ or $L$). The length- or temperature-dependence is therefore
		not a defect but the \emph{signature of a dynamic projection} (cf.\ Dok.~184,
		energy hierarchy $E_N = \hbar c / L_N$, and Landauer bound $k_B T \ln 2$).
	\end{insight}

	In the picture of the static/dynamic demarcation (Dok.~261) this is the
	\emph{layer-2} answer: an energy per bit exists only once the static structure is
	embedded in the time-flow and projected onto an absolute scale. This is where
	$\xi_{\text{UIFT}}$ lives -- temperature-bound, CMB-referenced, not universal.

	The static answer is different. The fundamental unit of information in the static
	context is not an energy but the \emph{winding number} on the compact torus:

	\begin{derivation}
		\textbf{The static unit of information: the winding number.}
		On the compact directions of the torus, winding numbers are topologically
		conserved and integer-valued (Dok.~181). A winding number is
		\begin{itemize}
			\item \textbf{dimensionless and countable} -- an integer, not an energy value;
			\item \textbf{reference-free} -- independent of $L$ and $T$, because it lies
			      \emph{before} any projection;
			\item \textbf{universal} -- the same counting unit at every scale.
		\end{itemize}
		This is exactly the state-space size that, according to Dok.~184, remains
		untouched by the $\xi$-dynamics (``the $\xi$-geometry sets the dynamics, not the
		state-space size''): the bit as a \emph{state} is static; only its energetic
		\emph{realization} is dynamic.
	\end{derivation}

	The two earlier statements thus resolve without contradiction:
	\emph{information as energy} is the dynamic projection (layer~2) --
	scale-dependent, without a universal unit. \emph{Information as counting}
	(winding, bit) is the static unit (layer~1) -- dimensionless and universal. The
	conversion between them is the \emph{embedding price} (Dok.~185); it explains why
	a ``bit'' acquires an energy only through $k_B T$ or $\hbar c/L$ -- and why that
	energy necessarily depends on the scale.

	In this light the Vopson comparison appears as a projection ratio: Vopson's
	$\xi_{\text{UIFT}} = k_B T_{\text{CMB}} \ln 2 / (m_e c^2)$ is a layer-2 quantity
	(bound to the CMB temperature through $k_B T_{\text{CMB}}$), while
	$\xi_{\text{FFGFT}}$ is the layer-1 invariant. The factor $\approx 418\,521$ is
	nothing other than the projection factor between a dynamic and a static
	information quantity.

	\section{The Interwoven Network of Constants}
	
	\subsection{The Fundamental Formula Network}
	
	\begin{derivation}
		\textbf{SI constants are mathematically linked:}
		
		Since the SI reform 2019, all fundamental constants are connected by exact mathematical relationships:
		
		\begin{align}
			\alpha &= \frac{e^2}{4\pi\varepsilon_0\hbar c} \quad \text{(exact definition)} \\
			\varepsilon_0 &= \frac{e^2}{2\alpha h c} \quad \text{(derived from above)} \\
			\mu_0 &= \frac{2\alpha h}{e^2 c} \quad \text{(via } \varepsilon_0\mu_0c^2 = 1) \\
			k_B &= \frac{R}{N_A} \quad \text{(definition of Boltzmann constant)}
		\end{align}
	\end{derivation}
	
	\subsection{The Geometric Boundary Condition}
	
	\begin{insight}
		\textbf{T0 theory reveals why these specific values are geometrically necessary:}
		
		\begin{equation}
			\alpha = \xi \cdot E_0^2 = \frac{1}{137{,}036} \quad \text{(geometric derivation)}
		\end{equation}
		
		This fundamental relationship forces the specific numerical values of the interwoven constants:
		
		\begin{equation}
			\frac{e^2}{4\pi\varepsilon_0\hbar c} = \frac{1}{137{,}036} \quad \text{(geometric boundary condition)} (See \cite{resolving_constants_alpha_En}.)
		\end{equation}
	\end{insight}
	
	\section{The Nature of Physical Constants}
	
	\subsection{Translation Conventions vs. Physical Quantities}
	
	\begin{keyresult}
		\textbf{Constants fall into three categories:}
		\begin{enumerate}
			\item \textbf{The single fundamental parameter:} $\xi = \frac{4}{3} \times 10^{-4}$
			
			\item \textbf{Geometric quantities derivable from $\xi$ (absolute values with one mass or energy anchor):}
			\begin{itemize}
				\item Particle masses (electron, muon, tau, quarks) (See \cite{particle_masses_En}.)
				\item Coupling constants ($\alpha$, $\alpha_s$, $\alpha_w$)
				\item Gravitational constant $G$
				\item Planck length $l_P$
			\end{itemize}
			
			\item \textbf{Pure translation conventions (SI unit definitions):}
			\begin{itemize}
				\item $\hbar$ (defines energy-time relationship)
				\item $e$ (defines charge scale)
				\item $k_B$ (defines temperature-energy conversion)
				\item $c$ (fixed by the definition of the metre since 1983)
			\end{itemize}
		\end{enumerate}
	\end{keyresult}
	
	\begin{warning}
		\textbf{Critical clarification about the speed of light:}
		
		The speed of light occupies a unique position in this classification:
		
		\begin{itemize}
			\item \textbf{In natural units ($c = 1$):} $c$ is a mere convention establishing how we relate length and time
			
			\item \textbf{In SI units:} The numerical value $c = 299\,792\,458$ m/s has been fixed by the definition of the metre since 1983. The ratio
			\begin{equation}
				c^{\text{SI}} = \frac{l_P^{\text{SI}}}{t_P^{\text{SI}}} = \frac{1{,}616 \times 10^{-35} \text{ m}}{5{,}391 \times 10^{-44} \text{ s}} = 299\,792\,458 \text{ m/s}
			\end{equation}
			is an identity because $t_P = l_P/c$, not a determination of $c$ by $\xi$.
		\end{itemize}
		
		\textbf{The consequence:} In natural units $c = 1$ is a convention, in SI units a fixing of the metre; in neither case does $\xi$ determine the numerical value of $c$.
	\end{warning}
	
	\subsection{The SI Reform 2019: Geometric Calibration Realized}
	
	The 2019 redefinition fixed constants by definition:
	\begin{align}
		c &= 299\,792\,458 \text{ m/s} \\
		\hbar &= 1{,}054571817... \times 10^{-34} \text{ J}\cdot\text{s} \\
		e &= 1{,}602176634 \times 10^{-19} \text{ C} \\
		k_B &= 1{,}380649 \times 10^{-23} \text{ J/K}
	\end{align}
	
	\begin{insight}
		This fixation implements the unique calibration consistent with $\xi$-geometry. The apparent arbitrariness conceals geometric necessity.
	\end{insight}
	
	\section{The Mathematical Necessity}
	
	\subsection{Why Constants Must Have Their Specific Values}
	
	\begin{derivation}
		\textbf{The interlocked system:}
		
		Given the fixed values and their mathematical relationships:
		
		\begin{align}
			h &= 2\pi\hbar = 6{,}62607015 \times 10^{-34} \text{ J}\cdot\text{s} \\
			\alpha &= \frac{e^2}{4\pi\varepsilon_0\hbar c} = \frac{1}{137{,}035999084} \\
			\varepsilon_0 &= \frac{e^2}{2\alpha h c} = 8{,}8541878128 \times 10^{-12} \text{ F/m} \\
			\mu_0 &= \frac{2\alpha h}{e^2 c} = 1{,}25663706212 \times 10^{-6} \text{ N/A}^2
		\end{align}
		
		These are not independent choices but mathematically enforced relationships. (See \cite{mathematical_structure_En} for mathematical structure.)
	\end{derivation}
	
	\subsection{The Geometric Explanation}
	
	\begin{historical}
		\textbf{Sommerfeld's unwitting geometric calibration}
		
		Arnold Sommerfeld's calibration in 1916 to $\alpha \approx 1/137$ established the SI system on geometric foundations. T0 theory reveals this was not coincidental but reflected the fundamental value $\alpha = 1/137{,}036$ derived from $\xi$. (See \cite{137_En}.)
	\end{historical}
	
		\section{Conclusion: Geometric Unity}
	
	\begin{keyresult}
		\textbf{Complete parameter freedom achieved:}
		\begin{itemize}
			\item \textbf{Single input:} $\xi = \frac{4}{3} \times 10^{-4}$
			
			\item \textbf{Derivable from $\xi$ (absolute SI values with one mass or energy anchor):}
			\begin{itemize}
				\item \textbf{First:} Particle masses including the electron via the ladder $m_e = \frac{4}{3}\,\xi^{3/2}\,v$ (Doc.~006), with $v$ as the scale
				\item \textbf{Then:} Gravitational constant: $G = \xi^2/(4m_e) \times$ (conversion factors)
				\item \textbf{Then:} Planck length: $l_P = \sqrt{\hbar G/c^3}$ with the complete $G$
				\item \textbf{Also:} Characteristic T0 length: $L_0 = \xi \cdot l_P$ (spacetime granulation)
				\item Coupling constants: $\alpha$, $\alpha_s$, $\alpha_w$
			\end{itemize}
			
			\item \textbf{Translation conventions (not derived, define units):}
			\begin{itemize}
				\item $\hbar$ defines energy-time relation in SI units
				\item $e$ defines charge scale in SI units
				\item $k_B$ defines temperature-energy conversion (historical)
				\item $c$ defines the metre (fixed by definition since 1983)
			\end{itemize}
			
			\item \textbf{Mathematical necessity:} Constants intertwined through exact formulas
			
			\item \textbf{Geometric foundation:} SI 2019 unwittingly implements $\xi$-geometry
		\end{itemize}
	\end{keyresult}
	
	\begin{center}
		\fbox{\parbox{0.9\textwidth}{
				\textbf{Final insight:} The universe is pure geometry, encoded in $\xi$. The complete derivation chain is:
				
				$\xi \to \{m_e, m_\mu, m_\tau, ...\} \to G \to l_P$
				
				with $L_0 = \xi \cdot l_P$, expressing the fundamental sub-Planck scale of spacetime granulation.
				
				\textbf{The profound mystery resolved:} Why does the Planck length, derived from $\xi$-geometry and one anchor, agree with the Planck length calculated from experimentally measured $G$? Because \emph{both describe the same geometric reality}: form and order of magnitude follow from $\xi$ and the anchor; the precision to $0.002\,\%$ arises only with the residual in $C_{\text{conv}}$. The SI reform 2019 unwittingly calibrated human measurement units to the fundamental $\xi$-geometry of the universe.
				
				This is no coincidence -- it is geometric necessity. Only $\xi$ is fundamental; everything else either follows from geometry or defines how we measure this geometry.
		}}
	\end{center}
	
	
	
	\section*{Appendix: Complete Derivation Chain}
	
	\begin{derivation}
		\textbf{From geometric parameter to measurable quantities:}
		
		1. Basic parameter: $\xi = \frac{4}{3} \times 10^{-4}$
		
		2. Electron mass: $m_e = \frac{4}{3}\,\xi^{3/2}\,v$ (Doc.~006), with $v$ as mass or energy anchor
		
		3. Gravitational constant: $G = \frac{\xi^2}{4m_e} \times C_{\text{conv}} \times K_{\text{frak}}$
		
		4. Planck length: $l_P = \sqrt{\hbar G/c^3}$ with $G$ from step 3
		
		5. Planck time: $t_P = l_P/c = l_P$ (natural units)
		
		6. Speed of light: $c = l_P/t_P$ is an identity because of step~5; $c = 299\,792\,458$ m/s is fixed by the definition of the metre
		
		7. Fundamental length: $L_0 = \xi \cdot l_P = 2{,}155 \times 10^{-39}$ m
		
		8. Fine structure constant: $\alpha = \xi \cdot E_0^2 = 1/137{,}036$
		
		\textbf{Consistency check:}
		\begin{align}
			\Delta G &= \left|\frac{G_{\text{T0}} - G_{\text{SI}}}{G_{\text{SI}}}\right| \approx 0.003\% \\
			\Delta l_P &= \left|\frac{l_P^{\text{T0}} - l_P^{\text{SI}}}{l_P^{\text{SI}}}\right| \approx 0.002\% \\
			\Delta \alpha &= \left|\frac{\alpha_{\text{T0}} - \alpha_{\text{SI}}}{\alpha_{\text{SI}}}\right| \approx 0.0005\%
		\end{align}
		
		The agreement of $G$ and $l_P$ at this level arises only with the residual in $C_{\text{conv}}$ and does not count as an independent precision result.
		
	\end{derivation}
	
	\section*{Glossary}
	
	\begin{description}
		\item[$\xi$] Fundamental geometric parameter, $\frac{4}{3} \times 10^{-4}$
		\item[$S_{T0}$] Mass scaling factor, $1$ MeV/$c^2$ (by construction)
		\item[$L_0$] Fundamental T0 length, $\xi \cdot l_P = 2{,}155 \times 10^{-39}$ m
		\item[$E_0$] Fundamental energy scale, $\sqrt{m_e \cdot m_\mu/K_{\text{frak}}} = 7.398$~MeV (bare $\sqrt{m_e m_\mu} = 7.348$~MeV)
		\item[$r_0(E)$] Characteristic length for energy $E$, $2GE$
	\end{description}
	
	
	
	\section{Bibliography}
	
	\begin{thebibliography}{99}
		
		\bibitem{T0_xi_origin_En}
		009\_T0\_xi\_origin\_En.pdf, .
		
		\bibitem{T0_gravitational_constant_En}
		012\_T0\_gravitational\_constant\_En.pdf, .
		
		\bibitem{gravitational_constant_En}
		045\_gravitational\_constant\_En.pdf, .
		
		\bibitem{T0_particle_masses_En}
		006\_T0\_particle\_masses\_En.pdf, .
		
		\bibitem{natural_units_systematics_En}
		015\_natural\_units\_systematics\_En.pdf, .
		
		\bibitem{fractal_correction_derivation_En}
		133\_fractal\_correction\_derivation\_En.pdf, .
		
		\bibitem{T0_natural_SI_En}
		014\_T0\_natural\_SI\_En.pdf, .
		
		\bibitem{T0_energy_En}
		010\_T0\_energy\_En.pdf, .
		
		\bibitem{T0_fine_structure_En}
		011\_T0\_fine\_structure\_En.pdf, .
		
		\bibitem{parameter_derivation_En}
		041\_parameter\_derivation\_En.pdf, .
		
		\bibitem{fine_structure_constant_En}
		044\_fine\_structure\_constant\_En.pdf, .
		
		\bibitem{unit_conventions_c_speed_En}
		134\_unit\_conventions\_c\_speed\_En.pdf, .
		
		\bibitem{circularity_constants_En}
		101\_circularity\_constants\_En.pdf, .
		
		\bibitem{temperature_units_CMB_En}
		061\_temperature\_units\_CMB\_En.pdf, .
		
		\bibitem{particle_masses_En}
		046\_particle\_masses\_En.pdf, .
		
		\bibitem{mathematical_structure_En}
		070\_mathematical\_structure\_En.pdf, .
		
		\bibitem{resolving_constants_alpha_En}
		043\_resolving\_constants\_alpha\_En.pdf, .
		
		\bibitem{137_En}
		087\_En.pdf, .
		
	\end{thebibliography}
	